% Einheit 3 von 5 – Additionsverfahren (bank/lineare-gleichungssysteme/e3.jsonl, zusammenbau v0.1)
%% TODO Zweigzeile Teil 1: was hier gelernt wird (Satz in Schülersprache aus dem Einheitstitel) – Einheit 3
\einheitenkopf[e3]{Einheit 3 von 5 · Additionsverfahren}
\zweigzeile{ab Kl. 9 (am Gymnasium ab Kl. 8) · keine P10-Aufgabe}
\setcounter{aufgabe}{18}

%% TODO \verfahren{…}: Verfahrensüberschrift in Sachsprache fehlt in der Bank (2.3 b)
%% TODO Titel: Ich-kann-Satz fehlt in der Bank (Platzhalter: Kettenname) – Nr. 19
%% TODO Anweisung: ein Satz über der ersten Teilaufgabe fehlt in der Bank – Nr. 19
\begin{aufgabe}{Addition}
%% TODO \rechenplatz{n}: Schrittzahl des Grundfalls fehlt in der Bank (nur bei fester Schreibform, 2.3 b)
\begin{teile}
\teil Gleiche Vorzahl oder Gegenzahl bei $y$? I: $2x + 3y = 9$, II: $5x - 3y = 14$. Kreuze an.\\ \kreuz{gleiche Vorzahl}\\ \kreuz{Gegenzahlen}\\ \kreuz{weder noch}
\end{teile}
\begin{gleichungsraster}[2]
\gl{\text{I: $x + y = 8$, II: $3x - y = 4$ – Lösung? Addiere die Gleichungen.}} &
\gl{\text{I: $3x + 6y = 27$, II: $3x + y = 17$ – Lösung? Subtrahiere die Gleichungen.}} \\
\gl{\text{I: $2x + y = 7$, II: $3x - 2y = 7$ – Lösung? Vervielfache eine Gleichung und addiere oder subtrahiere.}} &
\gl{\text{I: $2x + 3y = 16$, II: $5x - 2y = 21$ – Lösung? Vervielfache beide Gleichungen und addiere oder subtrahiere.}} \\
\gl{\text{I: $-2x + 3y = 14$, II: $2x + y = 2$ – Lösung? Löse mit dem Additionsverfahren.}} &
\gl{\text{I: $0{,}5x + y = 4$, II: $1{,}5x - y = 4$ – Lösung? Löse mit dem Additionsverfahren.}} \\
\gl{\text{I: $2x + y = 5$, II: $4x + 2y = 7$ – wie viele Lösungen? Rechne mit dem Additionsverfahren.}} &
\end{gleichungsraster}
\begin{teile}
\teil I: $y = 3x - 2$, II: $y = -x + 6$ – welches Verfahren passt? Entscheide und begründe.
\end{teile}
\begin{gleichungsraster}[2]
\gl{\text{I: $3x + 4y = 23$, II: $2x - 8y = -6$ – Lösung? Löse mit dem Additionsverfahren, mache die Probe in beiden Gleichungen und begründe, warum das Additionsverfahren hier passt.}} & \\
\end{gleichungsraster}
\end{aufgabe}

%% TODO \verfahren{…}: Verfahrensüberschrift in Sachsprache fehlt in der Bank (2.3 b)
%% TODO Titel: Ich-kann-Satz fehlt in der Bank (Platzhalter: Kettenname) – Nr. 20
%% TODO Anweisung: ein Satz über der ersten Teilaufgabe fehlt in der Bank – Nr. 20
\begin{aufgabe}{Addition, Sek II}
%% TODO \rechenplatz{n}: Schrittzahl des Grundfalls fehlt in der Bank (nur bei fester Schreibform, 2.3 b)
\begin{teile}
\teil Vorzahlen bei $x$ vergleichen: I: $-7x + 2y = 3$, II: $7x + 5y = 4$. Gleiche Vorzahl oder Gegenzahl? Kreuze an.\\ \kreuz{gleiche Vorzahl}\\ \kreuz{Gegenzahlen}\\ \kreuz{weder noch}
\end{teile}
\begin{gleichungsraster}[2]
\gl{\text{I: $2x - 3y = 4$, II: $-2x + 3y = 1$ – wie viele Lösungen? Addiere die Gleichungen und entscheide.}} &
\end{gleichungsraster}
\begin{teile}
\teil I: $x - 2y = 3$, II: $3x - 6y = \,?$ – welche Zahl gehört rechts in II, damit das System unendlich viele Lösungen hat? \leerfeld
\teil I: $2x - y = 4$, II: $ax + 3y = b$ mit reellen Zahlen $a$ und $b$ – für welche Werte hat das System keine Lösung? Gib Werte an und begründe mit einem Additionsschritt. \hfill (Abitur 2021 GK) \\ a = \leerfeld , b = \leerfeld
\end{teile}
\end{aufgabe}

%% TODO Titel: Ich-kann-Satz fehlt in der Bank (Platzhalter: Kettenname) – Nr. 21
%% TODO Anweisung: ein Satz über der ersten Teilaufgabe fehlt in der Bank – Nr. 21
\begin{aufgabe}{Dezimalzahlen und Brüche}
\begin{gleichungsraster}[2]
\gl{\text{I: $\frac{1}{2}x + y = 5$, II: $x - y = 1$ – Lösung? Löse mit dem Additionsverfahren.}} & \\
\end{gleichungsraster}
\end{aufgabe}

%% TODO Titel: Ich-kann-Satz fehlt in der Bank (Platzhalter: Kettenname) – Nr. 22
%% TODO Anweisung: ein Satz über der ersten Teilaufgabe fehlt in der Bank – Nr. 22
\begin{aufgabe}{Addition – Fehler finden, Begründen}
\begin{teile}
\teil Emil löst I: $x + 2y = 7$, II: $3x - 4y = 1$: \rechnung{2 \cdot \text{I}: 2x + 4y &= 7 \\ \text{plus II}: 5x &= 8} Finde den Fehler.
\teil Ida löst I: $5x + 3y = 17$, II: $5x - y = 1$: \rechnung{\text{I} - \text{II}: 3y - y &= 16 \\ 2y &= 16 \\ y &= 8} Finde den Fehler.
\teil I: $x - 3y = 2$, II: $ax + 6y = 1$ mit einer reellen Zahl $a$. Lara soll $a$ so wählen, dass das System keine Lösung hat, und wählt $a = 2$. Finde den Fehler.
\teil Warum darf man die Gleichungen I und II eines Systems addieren? Begründe.
\teil Nach dem Addieren zweier Gleichungen steht $0 = 4$. Warum hat das System dann keine Lösung? Begründe.
\end{teile}
\end{aufgabe}
